% Chapitre 1 : la méthode Elenchus, expliquée simplement.
\xdef\labelcourant{ch:elenchus}
\chapter{La méthode Elenchus, expliquée simplement}

Si je devais expliquer à quelqu'un qui n'est pas scientifique ce que j'ai fait cette année, je ne commencerais pas par
les théorèmes. Je commencerais par une question que chacun peut se poser : \emph{comment savoir si une machine dit
vrai ?} Les agents
d'IA écrivent vite et bien. Ils écrivent aussi, avec la même assurance, des choses fausses. Ce chapitre raconte la
méthode que nous avons fini par adopter pour ne plus les croire sur parole. Elle n'a rien de magique. Elle tient en
quelques règles qu'une cuisinière, un juge ou un inspecteur des ponts connaissent déjà.

\section{Un nom grec pour une idée simple}

\emph{Elenchos}, en grec ancien, désigne l'examen par questions, celui que pratiquait Socrate. Socrate ne cherchait pas
à prouver qu'il avait raison. Il interrogeait son interlocuteur jusqu'à trouver l'endroit où ses affirmations se
contredisaient. Ce qui résistait à l'examen méritait d'être cru, au moins un temps ; le reste devait être abandonné.

Nous avons repris l'idée en la transformant en ingénierie. Au lieu d'un philosophe qui pose des questions, il y a des
\emph{règles} que chaque résultat doit traverser avant d'être publié, et des \emph{programmes} qui posent une partie des
questions à notre place. La méthode a été écrite en septembre 2026, dans un article qui décrit la façon de travailler
d'un petit laboratoire assisté par l'IA, du choix des questions jusqu'à la publication.\reference{« A laboratory method
for AI-assisted science », 18 septembre 2026. \doi{10.5281/zenodo.22829767}} Elle a ensuite été appliquée, de plus en plus
strictement, à toutes les contributions de ce livre.

Voici les règles, une par une.

\begin{encadre}[Les sept règles d'Elenchus, en une page]
\begin{enumerate}[leftmargin=1.6em]
\item \textbf{Écrire l'affirmation avant le test}, et dire à l'avance ce qui la ferait tomber.
\item \textbf{Donner un niveau à chaque affirmation} : A, L, C ou X, selon qui l'a vérifiée.
\item \textbf{Faire un test qui doit échouer et un test qui doit réussir}, pour vérifier que le test fonctionne.
\item \textbf{Séparer celui qui fabrique de celui qui vérifie.}
\item \textbf{Publier les résultats négatifs}, avec le même soin que les autres.
\item \textbf{Confier la logique à un juge impartial}, le noyau de Lean, et savoir comment on peut le tromper.
\item \textbf{Tenir un registre} de toutes les affirmations, avec leur niveau et leurs preuves, et le laisser dire
« non » quand il manque quelque chose.
\end{enumerate}
\end{encadre}

\section{Règle 1 : écrire l'affirmation avant le test}

Un archer malhonnête tire d'abord sa flèche, puis peint la cible autour. Il touche toujours le centre. En science, la
tentation est la même : on fait un calcul, on regarde le résultat, puis on décide après coup ce qu'on cherchait. Le
remède est connu depuis longtemps en médecine, où il s'appelle le \emph{pré-enregistrement} : avant de lancer
l'expérience, on écrit, daté et archivé, ce que l'on prédit, et ce qui, si on l'observe, prouvera qu'on avait tort.

Dans ce projet, chaque hypothèse importante a été pré-enregistrée dans un fichier, avant le calcul. Le fichier dit
l'affirmation, la probabilité que l'auteur lui donnait de survivre au test, et la « règle d'arrêt » : ce qui la tuerait.
Par exemple, la conjecture que nous appellerons QU-1 (elle revient plus loin dans ce chapitre) a été inscrite avec une
chance de survie estimée à 0,6, c'est-à-dire un peu plus d'une chance sur deux.

Pré-enregistrer ne garantit pas qu'on ne se trompe pas. Cela garantit qu'on ne peut pas \emph{se raconter} qu'on ne
s'est pas trompé. Et cela oblige à une forme d'honnêteté supplémentaire : quand le travail réel s'écarte de ce qui
avait été annoncé, il faut le dire. Plusieurs articles de ce livre ont une section intitulée « déviations par rapport
au pré-enregistrement ». L'un d'eux, au chapitre~\ref{ch:c07}, signale qu'une de ses propres sous-hypothèses
pré-enregistrées était tout simplement fausse.

\section{Règle 2 : chaque affirmation porte un niveau}

La deuxième règle est la plus visible dans ce livre : chaque affirmation porte une étiquette qui dit \emph{qui l'a
vérifiée}. Il y a quatre niveaux, et chacun a son image.

\begin{figure}[htbp]
\centering
\begin{tikzpicture}[
  niv/.style={rectangle,rounded corners=3pt,minimum width=5.3cm,minimum height=1.55cm,text width=4.9cm,
    align=left,inner sep=6pt,font=\small},
  lettre/.style={circle,minimum size=0.95cm,font=\sffamily\bfseries\large,text=white,inner sep=0pt}]
  \node[niv,fill=okvert!14!ivoire,draw=okvert,line width=0.8pt] (A) at (0,0)
    {\lettreniv{okvert}{A}\textbf{« Le juge a regardé. »}\\ Un programme, le noyau de Lean, a contrôlé chaque pas logique.};
  \node[niv,fill=okbleu!12!ivoire,draw=okbleu,line width=0.8pt] (L) at (5.75,0)
    {\lettreniv{okbleu}{L}\textbf{« C'est écrit dans un livre sérieux. »}\\ Un résultat publié, lu et cité précisément.};
  \node[niv,fill=okorange!14!ivoire,draw=okorange,line width=0.8pt] (C) at (0,-1.95)
    {\lettreniv{okorange}{C}\textbf{« Un raisonnement. »}\\ Mené ici par des agents d'IA ; personne ne l'a vérifié à la machine.};
  \node[niv,fill=okpourpre!14!ivoire,draw=okpourpre,line width=0.8pt] (X) at (5.75,-1.95)
    {\lettreniv{okpourpre}{X}\textbf{« Une esquisse au crayon. »}\\ Un calcul en nombres à virgule : il oriente, il ne prouve rien.};
\end{tikzpicture}
\caption{Les quatre niveaux de confiance, avec les mots de tous les jours utilisés dans tout le livre.}
\label{fig:niveaux}
\end{figure}

\paragraph{\niveau{A} « le juge a regardé »} Un programme, le \emph{noyau} de Lean, a contrôlé chacun des pas logiques
de la démonstration. Ce n'est pas l'ordinateur qui « pense » ; il vérifie qu'aucun pas ne saute. Pour mériter ce
niveau, trois contrôles doivent passer : la démonstration n'utilise que trois règles de base de la logique, toujours
les mêmes (leurs noms techniques sont \texttt{propext}, \texttt{Classical.choice} et \texttt{Quot.sound}) ; un
second programme a relu la démonstration de bout en bout, ce qu'on appelle le \emph{rejeu} ; et une recherche des
raccourcis qui permettent de contourner le juge n'a rien trouvé. Quand le rejeu manque, le livre écrit
\niveau{As}.

\paragraph{\niveau{L} « c'est écrit dans un livre sérieux »} Le résultat a été publié dans la littérature
scientifique, et nous l'avons lu et cité avec la page et la ligne. Nous n'avons pas refait la preuve, mais d'autres
l'ont relue avant nous.

\paragraph{\niveau{C} « un raisonnement, mené ici par des agents d'IA, que personne n'a vérifié à la machine »}
C'est un raisonnement écrit, souvent juste, que personne n'a soumis au juge. Dans les sciences, ce raisonnement est
d'habitude celui d'un chercheur ; dans ce livre, il a presque toujours été mené par un agent d'IA et relu par un autre
agent. C'est une raison de plus d'être prudent.

\paragraph{\niveau{X} « une esquisse au crayon »} C'est un calcul fait par ordinateur avec des nombres à virgule, comme
ceux d'une calculatrice. Ces nombres sont arrondis à chaque opération ; ils donnent une très bonne idée, et parfois une
idée fausse. Une simulation, même énorme, reste une esquisse.

\medskip
Une affirmation complète mélange souvent plusieurs niveaux. Dans ce cas, \emph{elle vaut ce que vaut son maillon le
plus faible}. Un exemple le montrera à la fin du chapitre.

\begin{analogie}
Pensez aux étiquettes d'un pot de confiture. « Fait maison » ne dit pas la même chose que « contrôlé par le service
d'hygiène », ni que « recette de grand-mère, jamais analysée ». Aucune étiquette n'est honteuse ; ce qui serait
malhonnête, c'est de coller « contrôlé » sur un pot que personne n'a ouvert. L'image s'arrête là : un pot contrôlé
peut encore être mauvais, et un théorème de niveau A ne vaut que ce que vaut son énoncé, et le juge lui-même peut, très
rarement, se tromper ; c'est la question de l'énoncé qui reste ouverte, comme on le verra.
\end{analogie}

\section{Règle 3 : un test qui doit échouer, et un test qui doit réussir}

Un détecteur de fumée qui ne sonne jamais est très rassurant, jusqu'au jour de l'incendie. C'est pourquoi on appuie
sur son bouton de test : on vérifie qu'il \emph{sait} sonner. Les scientifiques appellent cela des \emph{contrôles}.
Un contrôle positif est un cas où l'on sait que la réponse doit être « oui » ; un contrôle négatif, un cas où l'on sait
qu'elle doit être « non ». Si le test répond « oui » aux deux, il ne vaut rien.

Dans les travaux de l'automne, aucune mesure n'est rapportée sans ses deux contrôles. Quelques exemples que l'on retrouvera dans les
chapitres :
\begin{itemize}
\item Pour vérifier que le juge Lean refuse bien une démonstration inachevée, on lui soumet exprès une démonstration
\emph{trouée}, avec la mention « à finir » (en Lean, ce mot s'écrit \texttt{sorry}). Le contrôle réussit quand le trou
est détecté.
\item Pour vérifier un certificat de calcul, on lui donne aussi un cas où il \emph{doit} échouer. Dans l'exemple de la
fin du chapitre, la même méthode, appliquée à un autre paramètre où le résultat attendu est différent, « échoue, comme
elle le doit ».
\item Pour vérifier un auditeur, on cache des erreurs dans le travail qu'on lui confie. Au chapitre~\ref{ch:c07},
l'auditeur a planté des erreurs dans une copie des fichiers, et elles ont été attrapées. Au chapitre~\ref{ch:c13}, une
copie volontairement abîmée d'un grand certificat a été rejetée.
\end{itemize}
Même la couverture de ce livre a ses contrôles : le programme qui la dessine vérifie que chaque cellule a six côtés et
une aire de~1, puis vérifie que ce test échoue sur un réseau légèrement déformé.

\begin{analogie}[Une image : le détecteur de fumée]
Un test qui n'a jamais échoué n'a jamais été testé. Avant de croire un « tout va bien », on vérifie que l'appareil
sait dire « attention ». L'image a une limite : un détecteur qui sonne au test peut encore manquer un feu d'un autre
genre. Un contrôle montre que le test marche sur \emph{ce} cas, pas sur tous.
\end{analogie}

\section{Règle 4 : celui qui fabrique n'est pas celui qui vérifie}

Le cuisinier qui goûte sa propre sauce la trouve toujours à son goût : il sait ce qu'il a voulu y mettre. C'est pour
cela qu'on fait goûter quelqu'un d'autre, qui ne connaît pas la recette et qui n'a aucune raison de faire plaisir au
cuisinier.

Dans les travaux de la fin de l'année, l'agent d'IA qui produit un résultat n'est pas celui qui le vérifie. Un second agent, lancé à part,
reçoit les fichiers, les recompile, refait les calculs avec ses propres programmes et rend un verdict écrit dans un
fichier d'audit. Les verdicts ont des noms précis : « accepté avec corrections », « confirmé avec corrections ». Puis
l'auteur applique les corrections, et un second fichier, le fichier de révision, en garde la trace. Par exemple,
l'audit des distances entre galaxies (chapitre~\ref{ch:c01}) a relevé treize défauts d'attribution et de portée,
corrigés avant la publication ; celui d'un des derniers articles de l'année, sur les forces écrantées, a conduit à douze
révisions.\reference{\doi{10.5281/zenodo.23663551}}

Cette séparation a une limite évidente, et nous la disons : le goûteur est lui aussi une IA, souvent du même fabricant.
Il ne connaît pas la recette, mais il a peut-être les mêmes papilles. C'est une raison pour laquelle un audit par IA
n'est pas une relecture humaine (nous y revenons plus bas).

\begin{analogie}[Une image : la recette goûtée par un autre]
On ne demande pas au cuisinier si son plat est bon ; on le fait goûter à quelqu'un qui n'a pas cuisiné, et on note ce
qu'il dit. L'image s'arrête là où commence la science : ici, le goûteur ne donne pas seulement son avis, il refait la
recette de son côté et compare les deux plats, gramme par gramme.
\end{analogie}

\section{Règle 5 : on publie les résultats négatifs}

Un résultat négatif, c'est une idée qui ne marche pas : une formule que les données contredisent, une conjecture
réfutée, une méthode qui ne s'étend pas. La tentation est grande de les ranger dans un tiroir. Mais un tiroir plein
d'échecs cachés fausse l'image que tout le monde se fait d'un sujet : on ne voit que ce qui a réussi.

Ce livre fait l'inverse. Trois des vingt contributions sont des « leçons » au sens strict, et beaucoup d'autres
contiennent une réfutation. Au chapitre~\ref{ch:c05}, les mesures de l'hélium superfluide contredisent une formule que
le projet avait lui-même proposée quelques semaines plus tôt. Aux chapitres~\ref{ch:c12} et~\ref{ch:c11}, une belle
idée (« le réseau triangulaire est aussi le meilleur dans le monde quantique ») est démentie, d'abord par le calcul,
puis avec un certificat. Au
chapitre~\ref{ch:c20}, une affirmation du lot du 19 octobre est retirée parce que la mesure la contredit.

Un résultat négatif bien fait n'est pas un échec : c'est une carte qui dit « ne passez pas par là ». Il apprend
souvent plus que le succès qu'on espérait.

\section{Règle 6 : un juge impartial, et comment on peut le tromper}

\subsection{Le juge}

Lean est un langage dans lequel on écrit des démonstrations mathématiques si détaillées qu'un ordinateur peut les
contrôler. Au centre de Lean se trouve un petit programme, le \emph{noyau}, dont la seule tâche est de vérifier
chaque pas : celui-ci découle-t-il vraiment du précédent, selon les règles admises ? Le noyau ne sait pas inventer une
démonstration. Il sait seulement dire « oui, chaque pas est correct » ou « non, ici quelque chose ne va pas ».

C'est un correcteur qui ne fatigue pas, ne flatte pas et se trompe rarement. Il ne fatigue pas : il relit la centième
page avec la même attention que la première ; l'audit du chapitre~\ref{ch:c15} lui a fait relire 88\,400
déclarations. Il ne flatte pas : il ne sait pas qui a écrit la démonstration, une IA célèbre ou un débutant, et cela
lui est égal. Il se trompe rarement : son code est petit et public, conçu pour pouvoir être relu ; mais rarement ne
veut pas dire jamais.

\begin{analogie}[Une image : le procès]
Imaginez un procès. L'agent d'IA qui a produit le résultat est la défense : il plaide que c'est vrai. L'agent qui
l'audite est l'accusation : il cherche la faille. Le noyau de Lean est le juge, et ce juge a une particularité : il ne
lit que les pièces écrites et la loi, jamais les plaidoiries. Les beaux discours ne l'impressionnent pas. Mais le juge
ne décide pas de tout. Il ne dit pas si le procès portait sur la bonne question : c'est à l'humain, à la fin, de
décider ce que le verdict veut dire et s'il mérite d'être publié. L'image s'arrête là : un vrai juge pèse des
témoignages, alors que le noyau ne pèse rien, il vérifie.
\end{analogie}

\subsection{Quatre façons de tromper le juge, et comment on s'en aperçoit}

Le juge est impartial, mais il n'est pas à l'abri d'une ruse. L'année nous en a appris quatre, et chacune a donné une
règle.

\paragraph{Le devoir rendu avec une page « à finir »} En Lean, on peut écrire \texttt{sorry} à la place d'un morceau de
démonstration : c'est une façon de dire « je reviendrai plus tard ». Le programme accepte alors de compiler le fichier,
et il \emph{termine sans erreur}. Pendant l'été, un article affirmait que ses vérifications Lean étaient réussies
parce que le programme avait terminé sans erreur. Or un \texttt{sorry} termine, lui aussi, sans erreur. Le remède : ne
jamais se fier au simple « terminé sans erreur », mais demander au juge la liste de ce sur quoi repose chaque
théorème. Si un trou s'y cache, il apparaît sous le nom \texttt{sorryAx}.

\paragraph{La règle ajoutée en douce} En Lean, on peut aussi déclarer un \emph{axiome} : une règle qu'on admet sans
démonstration. Si l'on admet une règle fausse, on peut tout démontrer, y compris que $1=2$. Un des travaux de ce
livre l'a montré exprès (chapitre~\ref{ch:c18}) : un fichier « démontrant » $1=2$ grâce à un axiome caché se compile
sans la moindre erreur, et seul l'audit des axiomes le démasque. D'où la règle : la liste des axiomes de chaque
théorème doit se limiter aux trois règles de base citées plus haut.

\paragraph{La pièce glissée dans le dossier sans passer devant le juge} C'est la ruse la plus subtile, et elle a été
trouvée par nos propres audits. Lean permet, par une option réservée aux développeurs, d'ajouter une déclaration
\emph{sans la faire contrôler} par le noyau. Lors de l'audit des preuves d'OpenAI sur les équations des fluides
(chapitre~\ref{ch:c15}), l'auditeur a fabriqué exprès une telle pièce, qui « démontre » une absurdité, et a constaté
un résultat négatif qu'il a placé en tête de son rapport : \emph{la liste des axiomes ne détecte pas cette
contrefaçon}.\reference{Audit A2, 26 octobre 2026, \nolinkurl{AXIOMS.md} et \nolinkurl{SOUNDNESS_SCAN.md}.} Elle est propre ;
le juge n'a simplement jamais vu la pièce.

La parade, c'est le \emph{rejeu} : un second programme, \texttt{leanchecker}, relit depuis le début toutes les
déclarations enregistrées et les soumet de nouveau au noyau. Le 26 octobre, en même temps que le rejeu de nos propres modules, nous
avons refait l'expérience sur une contrefaçon de contrôle. La liste des axiomes répondait que la pièce « ne dépend
d'aucun axiome » : le contrôle habituel l'aurait laissée passer. Le rejeu, lui, s'est arrêté en erreur, avec le
message « le type de la déclaration ne correspond pas ».\reference{Rejeu du 26 octobre 2026, \nolinkurl{results/kernel_replay_2026-10-26/},
contrôle \texttt{KRForged}.} Le rejeu a été lancé sur 93 de nos modules de démonstration, qui ont tous passé ; une
recherche des quatre options de contournement connues dans 871 fichiers n'a rien trouvé.

\begin{analogie}[Une image : l'inspection du pont]
Un pont porte une plaque : « inspecté et conforme ». La plaque, c'est la liste des axiomes : elle dit ce que le
dossier contient. Mais un dossier peut contenir une page que l'inspecteur n'a jamais signée. Le rejeu, c'est
l'inspecteur qui revient, et qui refait le tour du pont boulon par boulon, sans regarder la plaque. L'image a une
limite : un pont s'use, une démonstration non. Une fois le rejeu passé, il n'y a pas à recommencer, sauf si l'on change
les fichiers ou les outils.
\end{analogie}

\paragraph{La bonne preuve d'une mauvaise question} La dernière ruse n'en est pas une, et c'est la plus importante. Le
juge vérifie que la démonstration prouve l'énoncé écrit. Il ne vérifie pas que l'énoncé écrit dit bien ce que nous
\emph{racontons} qu'il dit. Si j'écris en Lean « tout carré est un rectangle » en me trompant de définition, le juge
validera avec joie un théorème sans intérêt. Seul un lecteur qui compare l'énoncé formel avec la question posée peut
attraper ce genre d'erreur. Cette étape, que nous appelons \emph{l'audit de l'énoncé}, n'a encore été faite par aucun
humain pour les résultats de ce livre. Le registre, règle suivante, le dit à chaque fois.

\section{Règle 7 : le registre des affirmations}

Chaque article de la fin de l'année est accompagné d'un \emph{registre} : un fichier qui contient une ligne par
affirmation. Chaque ligne porte un identifiant, le niveau, l'énoncé en mots, une empreinte numérique de la preuve
(une sorte de code-barre qui change si l'on modifie le moindre fichier), et une case « audit humain ». Voici, traduite,
une vraie ligne du registre de l'audit des preuves sur les fluides (chapitre~\ref{ch:c15}) :

\begin{encadre}[Une ligne du registre]
\small
\textsf{\textbf{P2NS-A-0001}} \quad \niveau{A} \quad \textsf{audit humain : aucun}\par\smallskip
\emph{Énoncé :} le théorème Lean \texttt{navier\_stokes\_breakdown\_R3} affirme que, pour toute viscosité positive, il
existe une vitesse initiale et une force, lisses et décroissant très vite, pour lesquelles aucune solution lisse
d'énergie bornée des équations de Navier--Stokes n'existe. Notre propre audit des axiomes donne exactement les trois
règles de base ; le module a été reconstruit à partir des sources, et ses 53\,752 déclarations ont été rejouées par le
noyau. \emph{Un énoncé sur le terme Lean uniquement.}\par\smallskip
\emph{Preuve :} empreinte \texttt{sha256:0c86f9d7\dots}
\end{encadre}

Un programme, la \emph{porte du registre}, lit toutes les lignes et refuse de dire « tout va bien » tant qu'il manque
quelque chose. Pour cet article, sa réponse est : quatre fois le signal \texttt{LEDGER\_UNAUDITED\_TIER\_A}, « le noyau a
vérifié la preuve, et aucune personne n'a encore certifié que l'énoncé veut dire ce que dit sa glose », puis le code de
sortie~1, qui veut dire « non ».

C'est voulu. Presque tous les registres de ce livre se terminent ainsi, parce que l'audit humain des énoncés n'a pas
été fait. L'auteur a choisi de publier quand même, et de l'écrire dans chaque article. C'est déjà le cas de
l'article du chapitre~\ref{ch:c03}, dont le registre compte vingt-deux affirmations. Publier avec le drapeau rouge
levé, en le montrant, est plus honnête que de baisser le drapeau.

\section{Ce qu'un audit par IA n'est pas}

Il faut le dire clairement, parce que l'année nous a pris en défaut sur ce point : \textbf{un audit par un agent d'IA
n'est pas une relecture par des pairs}. La relecture par des pairs, dans la science, c'est l'examen d'un travail par
des chercheurs humains du domaine, qui engagent leur jugement et leur réputation. Rien dans ce livre n'a reçu cet
examen.

Les 19 et 20 octobre, des articles sont parus avec la mention « relu en cinq tours par des pairs ». Le 21 octobre,
l'un d'eux a publié une seconde édition qui retire cette mention : le fichier de relecture joint « a été produit par un
script qui n'appelle aucun relecteur ».\reference{Retrait du 21 octobre 2026 : \doi{10.5281/zenodo.23525655}} Dès le
19 octobre, un autre dépôt avait retiré un « rapport de relecture » écrit par un programme, avec des relecteurs
imaginaires.\reference{\doi{10.5281/zenodo.23468458}} Un mois plus tôt, l'audit de l'atelier lui-même avait trouvé un
script de « relecture par les pairs » qui n'appelait jamais de modèle.

D'où deux règles, que nous appliquons désormais à la lettre. Un fichier de relecture doit montrer qu'un relecteur a
réellement été appelé : qui, quand, avec quelle réponse. Et l'on écrit « audité par un agent d'IA », jamais « relu par
des pairs », tant qu'aucun humain du domaine ne l'a fait.

\section{Un exemple : une idée qui monte et descend les niveaux}

Voici, pour finir, l'histoire d'une seule affirmation, du brouillon au registre. Elle est racontée en détail aux
chapitres~\ref{ch:c12} et~\ref{ch:c11} ; ici, on ne suit que ses niveaux.

\paragraph{L'idée} On sait (c'est le fil de la deuxième partie) que le réseau triangulaire minimise l'énergie
\emph{immobile} de points qui se repoussent. Mais dans le monde quantique, les points ne sont jamais tout à fait
immobiles : même au zéro absolu, ils vibrent un peu. Cette vibration minimale a une énergie, l'\emph{énergie de point
zéro}. L'idée naturelle, baptisée QU-1, était : « le réseau triangulaire minimise aussi l'énergie de point zéro ». À ce
stade, c'est une conjecture : pas même un argument, une intuition. Elle est pré-enregistrée, avec une chance de survie
estimée à~0,6.

\paragraph{L'esquisse au crayon \niveau{X}} Avant d'essayer de la démontrer, on la teste par le calcul. Deux
formules, programmées de quatre façons, puis une cinquième fois par l'auditeur, calculent l'énergie de point zéro de nombreux réseaux. Résultat : pour huit des douze forces
testées, un autre réseau, stable, a une énergie de point zéro plus basse que le réseau triangulaire. QU-1 est réfutée,
mais seulement « au crayon » : des nombres à virgule, si concordants soient-ils, ne prouvent rien.\reference{Réfutation
numérique, 26 octobre 2026 : \doi{10.5281/zenodo.23653833}}

\paragraph{Le certificat \niveau{A} pour la logique, \niveau{C} pour l'ensemble} Pour une des forces, la cloche
gaussienne, on veut une preuve. On déforme très légèrement le réseau triangulaire, par un petit \emph{cisaillement},
comme on pousse le dessus d'une pile de cartes. Un calcul rigoureux par encadrements, qui garantit à chaque étape que
la vraie valeur est dans un intervalle (cela s'appelle l'arithmétique de boules), montre que l'énergie de point zéro
\emph{baisse} pour tout cisaillement assez petit : le réseau triangulaire n'est donc même pas un minimum local. Le
calcul a découpé le problème en 206\,468 cellules et aucune n'a échoué. Le noyau de Lean vérifie la logique de
l'argument et l'arithmétique exacte du certificat.\reference{Certificat, 26 octobre 2026 :
\doi{10.5281/zenodo.23655392}}

Mais les encadrements eux-mêmes ont été calculés par un programme \emph{extérieur} au noyau, et certaines étapes qui
relient les morceaux sont des raisonnements écrits par des agents d'IA, que ni une machine ni un humain n'a vérifiés. Alors l'article écrit, honnêtement : « le résultat principal est
de niveau C dans son ensemble ; seules sa logique et son arithmétique sont de niveau A ». C'est la règle du maillon le
plus faible.

\paragraph{Les contrôles} La même méthode, appliquée à une cloche plus étroite pour laquelle les calculs « au crayon »
ne prévoient pas de baisse, ne parvient pas à certifier de baisse, « comme elle le doit ». Des démonstrations
volontairement trouées, un axiome planté et un certificat falsifié ont été soumis au juge, qui les a tous refusés.

\paragraph{L'audit et le registre} Un agent distinct a refait les calculs ; un article suivant a étendu le résultat aux
forces écrantées, avec un audit et douze révisions.\reference{Forces écrantées, 26 octobre 2026 :
\doi{10.5281/zenodo.23663551}} Et la porte du registre répond toujours
« non » : aucune personne n'a encore certifié que les énoncés Lean disent bien ce que disent leurs gloses.

\begin{encadre}[Ce que l'exemple montre]
En quelques semaines, une même idée a été une intuition, une hypothèse pré-enregistrée, une réfutation au crayon, puis
un certificat mixte dont chaque morceau porte son niveau. À aucun moment on n'a écrit « démontré » pour ce qui ne
l'était pas. C'est cela, Elenchus : non pas avoir toujours raison, mais savoir à chaque instant \emph{à quel point} on
a raison, et le dire.
\end{encadre}

\section{Même les plus grands : vérifier la preuve d'une autre IA}

Les règles ne s'appliquent pas seulement à nos propres agents. Elles s'appliquent, avec la même rigueur, aux travaux
des grands laboratoires. En 2026, OpenAI a publié des démonstrations produites par une IA et vérifiées par
ordinateur. Il aurait été tentant de les prendre telles quelles : elles venaient d'une équipe réputée, et le juge
Lean les avait acceptées. Nous les avons pourtant examinées comme les nôtres, avant de bâtir dessus.

\paragraph{Une affirmation sur les charges électriques} La publication affirme aussi démontrer la conjecture de
Sandier et Serfaty, sur la disposition la plus économe de charges dans un fond neutralisant, par une réduction et un
très gros certificat de calcul. Un agent a refait tourner ce certificat : il s'est reproduit en environ huit minutes.
Il a vérifié qu'une copie volontairement abîmée est rejetée (c'est le contrôle négatif), et il a lu la réduction pas à
pas sans y trouver de trou. Le chapitre~\ref{ch:c13} le raconte, et son titre le dit : \emph{aucune preuve n'est
revendiquée}. Ni démontrée, ni vérifiée, ni réfutée par nous ; reproduite et auditée, ce qui est différent. Au passage,
l'audit a trouvé et réparé un trou dans \emph{notre} outil de certificat. Puis quatre des douze étapes du certificat
ont été refaites à l'intérieur du noyau (chapitre~\ref{ch:c14}) ; les huit autres attendent.

\paragraph{Des théorèmes sur les fluides} OpenAI a aussi publié, en Lean, des théorèmes d'« explosion » pour les
équations de Navier--Stokes et d'Euler, qui décrivent l'écoulement des fluides. Un agent les a reconstruits de façon
indépendante : 2\,440 modules recompilés, 88\,400 déclarations rejouées par le noyau, aucune échappatoire trouvée, et
les énoncés comparés aux textes officiels des problèmes (chapitre~\ref{ch:c15}). Ce travail ne prétend pas qu'un
problème du millénaire soit résolu, et le mérite des preuves revient à leurs auteurs. Il dit seulement, de façon
vérifiable par quiconque, ce qui a été contrôlé et ce qui ne l'a pas été.

\begin{analogie}[Une image : le contrôle technique]
Une voiture neuve d'une grande marque passe quand même le contrôle technique. Ce n'est pas un soupçon contre la
marque ; c'est la règle, la même pour tous. L'image s'arrête là : le contrôleur d'une voiture connaît la liste des
points à vérifier, alors que pour une preuve nouvelle, une partie du travail consiste à inventer cette liste.
\end{analogie}

\section{Les niveaux des vingt, d'un coup d'œil}

Le tableau de la fin du livre donne, pour chaque contribution, sa note de rigueur dans la grille de sélection, de 0 à
3. Seize obtiennent 3 : une preuve vérifiée par le juge et rejouée, ou un certificat de calcul recontrôlé par un
programme différent. Une obtient 2 (vérifiée par le juge, sans rejeu), deux obtiennent 1, et une obtient 0 :
l'instrument de calcul rusty-SUNDIALS, dont tous les nombres sont des nombres à virgule. C'est voulu : un instrument
n'a pas à être un théorème.

Il faut lire ce décompte avec prudence. La note est celle du \emph{meilleur} morceau de chaque travail, pas de son
ensemble ; elle n'est pas non plus tout à fait un niveau Elenchus, puisqu'un certificat recontrôlé sans rejeu du juge
y vaut autant qu'un rejeu. Les pastilles de niveau de chaque chapitre font foi. Une contribution de niveau A contient presque toujours aussi des simulations (X), des citations (L) et
des raisonnements (C), et sa conclusion principale vaut ce que vaut son maillon le plus faible, comme dans l'exemple
de QU-1. C'est pourquoi chaque chapitre a une section « Comment on s'est assuré que ce n'est pas faux », qui détaille
les niveaux morceau par morceau, et une section « Ce que ça ne dit pas ».

\section{Les règles appliquées à ce livre}

Ce livre essaie de suivre ses propres règles. Ses sources sont fixées à l'avance : un nombre qui ne figure pas dans le
dossier d'un travail n'y est pas utilisé. Les vingt contributions ont été choisies par une grille écrite et gelée avant
la moindre note, appliquée de la même façon à tous les candidats (c'est l'équivalent du pré-enregistrement). Chaque
chapitre suit le même plan, quel que soit le domaine, pour qu'aucun résultat ne soit mieux présenté qu'un autre
simplement parce qu'il est plus séduisant. Les erreurs et les retraits ont leur propre chapitre à la fin. Et ce livre,
comme ses résultats, n'a pas été relu par un scientifique humain extérieur : il porte donc, lui aussi, un drapeau
levé.

\section{Ce que la méthode ne fait pas}

Elenchus n'est pas une garantie. Elle ne remplace pas la relecture humaine ; elle ne dit pas si une question vaut la
peine d'être posée ; elle n'empêche pas de se tromper, elle rend les erreurs plus faciles à trouver. Et elle ne
s'applique pas toute seule : la rechute du 19 octobre est arrivée un mois \emph{après} que la méthode a été écrite.
Une règle sur le papier ne protège de rien. Ce qui protège, c'est de la faire tourner, à chaque fois, y compris quand
on est pressé et content de soi.

Les vingt chapitres qui suivent en sont l'application. Lisez le niveau avant la phrase.
